% !TeX root = RJwrapper.tex
\title{grip: Graph Drawing with Intelligent Placement in R}


\author{by Pawel Gajer, Jacques Ravel}

\maketitle

\abstract{%
The \textbf{grip} package brings the GRIP multiscale force-directed graph layout algorithm to R, providing scalable 2D and 3D layouts for graphs ranging from tens to thousands of vertices. The package offers a clean API built around a single core function (\texttt{grip.layout()}), validated parameter presets for common graph families, and---most distinctively---a principled layout quality scoring framework (\texttt{grip.score.layout()}, \texttt{grip.compare.layouts()}) that supports systematic parameter selection for real-world graphs where no ground-truth embedding exists. The C++ core is accessed via Rcpp, and the package includes graph generators, a plotting helper with 3D projection support, and a bundled 1,828-vertex microbial network from the Human Microbiome Project for realistic worked examples. This article describes the algorithm, the package design, the quality metrics, and demonstrates the full workflow on synthetic and real-data examples.
}

\section{Introduction}\label{introduction}

Visualizing graphs (networks) is a fundamental task in data analysis across
disciplines---from social network analysis to systems biology to knowledge graph
exploration. Force-directed layout algorithms, which model vertices as charged
particles and edges as springs, are the most widely used family of methods for
producing readable graph drawings. Classical algorithms such as
Fruchterman--Reingold \citep{fruchterman1991} and Kamada--Kawai \citep{kamada1989} produce
good layouts for small graphs but scale poorly: their \(O(n^2)\) or \(O(n^3)\)
per-iteration cost becomes prohibitive beyond a few hundred vertices.

For the large graphs routinely encountered in biological networks, social media
analysis, and knowledge graph construction, multiscale methods are needed. The
GRIP algorithm \citep{gajer2002, gajer2004} addresses this by building a hierarchy
of progressively coarser vertex subsets using maximal independent set (MIS)
filtrations, placing vertices coarse-to-fine, and refining with local
force-directed iterations at each level. This yields near-linear time complexity
in practice and produces high-quality layouts even for graphs with thousands of
vertices.

The \textbf{grip} package implements the GRIP algorithm in R with a C++ core
accessed via Rcpp \citep{eddelbuettel2011}. Beyond the layout computation itself,
the package makes three additional contributions. First, it ships validated
parameter presets for common graph families (meshes, trees, fractal carpets,
tori), so that users can obtain good layouts without manual tuning. Second, it
provides a principled layout quality scoring framework with metrics including
sampled stress, edge-length uniformity, non-edge separation, Procrustes
stability, and optional cluster separation when external labels are available.
Third, it offers a systematic comparison workflow
(\texttt{grip.compare.layouts()}) that evaluates multiple candidate
parameter settings across random seeds, enabling reproducible layout selection
for real-world graphs where no ground-truth embedding is known.

Several R packages provide graph layout capabilities. The \textbf{igraph} package
\citep{csardi2006} offers implementations of Fruchterman--Reingold,
Kamada--Kawai, and the distributed recursive layout (DrL) algorithm, among
others. The \textbf{ggraph} package \citep{pedersen2024} provides a grammar-of-graphics
layer for network visualization but delegates layout computation to igraph. The
\textbf{graphlayouts} package \citep{schoch2024} implements stress-based layouts with
pivot MDS initialization. The \textbf{visNetwork} package provides interactive
JavaScript-based visualization. The \textbf{grip} package is complementary to these
tools: it focuses specifically on the layout computation and quality assessment
problem, and its scoring functions can evaluate layouts computed by any package.

The author of this package (P. Gajer) is one of the original developers of the
GRIP algorithm. This R implementation was motivated by the need for a
well-tested, accessible version of GRIP within the R ecosystem, with the
added tooling for quality assessment that was not part of the original
publications.

\section{Algorithm overview}\label{algorithm-overview}

The GRIP algorithm produces graph layouts through a coarse-to-fine multiscale
process. This section provides an intuitive overview; full algorithmic details
are given in \citet{gajer2002} and \citet{gajer2004}.

\subsection{Maximal independent set filtration}\label{maximal-independent-set-filtration}

The algorithm begins by constructing a hierarchy of vertex subsets
\(V = V_0 \supset V_1 \supset \cdots \supset V_k\), where each \(V_{i+1}\) is a
maximal independent set (MIS) of the subgraph induced by \(V_i\) in the
graph-distance metric. The coarsest level \(V_k\) contains only a small number
of vertices (controlled by the \texttt{num\_init} parameter), while the finest level
\(V_0\) is the full vertex set. This filtration ensures that vertices at each
level are well-separated in the graph metric, providing a natural hierarchy
for multiscale placement.

\subsection{Coarse-to-fine placement}\label{coarse-to-fine-placement}

Layout computation proceeds from the coarsest level to the finest. The
vertices in \(V_k\) are placed first using an initial arrangement (either
random barycenter placement or circular placement in 2D, controlled by the
\texttt{placement} parameter). Then, for each successive level \(V_i\) (\(i = k-1,
\ldots, 0\)), the newly introduced vertices are positioned based on their
graph distances to already-placed neighbors using a local minimization of a
distance-based energy function.

\subsection{Local force-directed refinement}\label{local-force-directed-refinement}

After each level's vertices are placed, the algorithm performs local
force-directed refinement iterations. At each iteration, every vertex
experiences attractive forces from its graph neighbors (proportional to the
discrepancy between the current Euclidean distance and the target graph
distance) and optional repulsive forces from nearby non-neighbors. The
displacement of each vertex is modulated by an adaptive local temperature
that tracks whether the vertex is oscillating or moving consistently. The
temperature update rule uses two parameters: \texttt{r} controls the cooling rate,
and \texttt{s} provides a directional momentum boost when successive displacements
are aligned. The number of refinement iterations at intermediate levels is
set by \texttt{rounds}, while the final level uses \texttt{final\_rounds} iterations for
the most detailed refinement.

The \texttt{repulsion\_factor} parameter controls the strength of the repulsive term
at the finest level. Setting it to zero disables repulsion entirely, which
can be appropriate for tree-like graphs where repulsion tends to distort the
natural branching structure.

The \texttt{num\_nbrs} parameter controls how many graph-distance neighbors are
retained for the local force computation at each level. Larger values produce
more globally faithful layouts at the cost of additional computation.

\subsection{Computational complexity}\label{computational-complexity}

The MIS filtration construction is linear in the number of edges. At each
level, the force computation involves only the local neighborhood of each
vertex (bounded by \texttt{num\_nbrs}), so the per-level cost is \(O(n \cdot
\texttt{num\_nbrs})\). With \(O(\log n)\) levels in typical graphs, the overall
complexity is \(O(n \cdot \texttt{num\_nbrs} \cdot \texttt{rounds} \cdot \log
n)\), which is near-linear in practice for bounded-degree graphs.

\subsection{Multiscale trace}\label{multiscale-trace}

The \texttt{grip.layout.trace()} function captures snapshots of the layout at
intermediate stages of the multiscale process, which is useful for
understanding how the algorithm operates and for creating animated
visualizations. Figure @ref(fig:trace) shows four stages of a mesh layout,
from the initial coarse placement through progressive refinement to the
final layout.

\begin{verbatim}
library(grip)

edges <- edges.mesh(6, 6)
tr <- grip.layout.trace(
  edges, n = 36, dim = 2,
  preset = "mesh", seed = 31,
  trace = "level"
)

frame.idx <- unique(c(
  1L,
  max(2L, floor(length(tr$frames) / 3)),
  max(3L, floor(2 * length(tr$frames) / 3)),
  length(tr$frames)
))

op <- par(mfrow = c(1, 4), mar = c(1, 1, 2.5, 1))
labels <- c("Coarsest start", "Early refinement",
            "Late refinement", "Final layout")
for (i in seq_along(frame.idx)) {
  coords <- tr$frames[[frame.idx[i]]]
  active <- complete.cases(coords[, 1:2])
  xy <- coords[active, 1:2, drop = FALSE]
  active.edges <- edges[active[edges[, 1]] & active[edges[, 2]], , drop = FALSE]

  xlim <- range(xy[, 1]); ylim <- range(xy[, 2])
  xpad <- max(0.08 * diff(xlim), 0.2)
  ypad <- max(0.08 * diff(ylim), 0.2)

  plot(xy[, 1], xy[, 2], type = "n", asp = 1, axes = FALSE,
       xlab = "", ylab = "",
       xlim = xlim + c(-xpad, xpad),
       ylim = ylim + c(-ypad, ypad),
       main = labels[i])
  if (nrow(active.edges) > 0)
    segments(coords[active.edges[, 1], 1], coords[active.edges[, 1], 2],
             coords[active.edges[, 2], 1], coords[active.edges[, 2], 2],
             col = "gray80")
  points(xy[, 1], xy[, 2], pch = 16, cex = 0.6)
}
\end{verbatim}

\begin{figure}

{\centering \includegraphics[width=1\linewidth]{grip-paper_20260326_092635_files/figure-latex/trace-1} 

}

\caption{Multiscale trace of a 6x6 mesh layout using the mesh preset. From left to right: coarsest initialization with a handful of representative vertices, early refinement introducing additional vertices, late refinement improving spacing, and the final converged layout.}(\#fig:trace)
\end{figure}

\begin{verbatim}
par(op)
\end{verbatim}

\section{Package design and API}\label{package-design-and-api}

\subsection{Core layout function}\label{core-layout-function}

The primary function is \texttt{grip.layout()}, which accepts either an edge-list
matrix or an adjacency list with optional edge weights:

\begin{verbatim}
# Edge-list input: 8x8 mesh in 2D
edges <- edges.mesh(8, 8)
coords <- grip.layout(edges, n = 64, dim = 2, preset = "mesh", seed = 1)
\end{verbatim}

\begin{verbatim}
# Adjacency-list input with weights
adj_list <- list(c(2L), c(1L, 3L), c(2L, 4L), c(3L))
weight_list <- list(c(1.0), c(1.0, 2.0), c(2.0, 1.5), c(1.5))
coords <- grip.layout(adj_list = adj_list, weight_list = weight_list,
                       n = 4, dim = 2, seed = 12)
\end{verbatim}

The function returns an \(n \times d\) numeric matrix of coordinates, where
\(d \in \{2, 3\}\). Key parameters include \texttt{dim} (layout dimension), \texttt{placement}
(\texttt{"barycenter"} or \texttt{"circle"}), \texttt{rounds} and \texttt{final\_rounds} (refinement
iteration counts), \texttt{num\_init} (vertices at the coarsest level), \texttt{num\_nbrs}
(local neighborhood size), \texttt{r} and \texttt{s} (temperature adaptation), and
\texttt{repulsion\_factor} (finest-level repulsion strength). The \texttt{seed} argument
ensures reproducibility.

Disconnected graphs are handled automatically by default
(\texttt{disconnected\ =\ "components"}): each connected component is laid out
independently and the results are packed into a single coordinate matrix
with a gap-based arrangement.

\subsection{Presets}\label{presets}

Presets are validated parameter bundles for common graph families. Each preset
sets sensible defaults for the GRIP tuning parameters; any argument supplied
explicitly by the user overrides the corresponding preset value. The package
currently ships four presets:

\begin{longtable}[]{@{}
  >{\raggedright\arraybackslash}p{(\linewidth - 4\tabcolsep) * \real{0.3333}}
  >{\raggedright\arraybackslash}p{(\linewidth - 4\tabcolsep) * \real{0.3333}}
  >{\raggedright\arraybackslash}p{(\linewidth - 4\tabcolsep) * \real{0.3333}}@{}}
\toprule\noalign{}
\begin{minipage}[b]{\linewidth}\raggedright
Preset
\end{minipage} & \begin{minipage}[b]{\linewidth}\raggedright
Target family
\end{minipage} & \begin{minipage}[b]{\linewidth}\raggedright
Tuned on
\end{minipage} \\
\midrule\noalign{}
\endhead
\bottomrule\noalign{}
\endlastfoot
\texttt{"mesh"} & Rectangular grid/lattice & \(8 \times 8\) and \(12 \times 12\) meshes \\
\texttt{"carpet"} & Sierpinski carpet & Level 3 and 4 carpets \\
\texttt{"tree"} & Tree-like graph & Binary trees, depths 5--6 \\
\texttt{"torus"} & 3D torus/cylinder & Torus sizes \(8 \times 8\) through \(20 \times 20\) \\
\end{longtable}

For instance, the mesh preset sets \texttt{placement\ =\ "barycenter"},
\texttt{rounds\ =\ 128}, \texttt{final\_rounds\ =\ 128}, \texttt{num\_init\ =\ 12}, \texttt{num\_nbrs\ =\ 20},
\code{r = 0.10}, \texttt{s\ =\ 4.5}, and \texttt{repulsion\_factor\ =\ 1.5}. The tree preset
uses \texttt{placement\ =\ "circle"} in 2D and disables repulsion
(\texttt{repulsion\_factor\ =\ 0}), which is appropriate for branching structures.

\subsection{Plotting}\label{plotting}

The \texttt{grip.plot()} function provides a lightweight plotting helper. For 2D
layouts, it draws vertices and edges using base graphics. For 3D layouts, it
supports both interactive viewing via the optional \textbf{rgl} package
(\texttt{projection\ =\ "rgl"}, the default) and static orthographic projection
(\texttt{projection\ =\ "ortho"}) suitable for publications and vignettes.

\begin{verbatim}
edges <- edges.mesh(8, 8)
coords <- grip.layout(edges, n = 64, dim = 2, preset = "mesh", seed = 1)
grip.plot(coords, edges,
          pch = 16, cex = 0.6,
          axes = FALSE, xlab = "", ylab = "", frame.plot = FALSE,
          main = "8x8 mesh (mesh preset)")
\end{verbatim}

\begin{figure}

{\centering \includegraphics[width=1\linewidth]{grip-paper_20260326_092635_files/figure-latex/mesh-plot-1} 

}

\caption{An 8x8 mesh laid out with the mesh preset in 2D.}(\#fig:mesh-plot)
\end{figure}

The \texttt{grip.project.3d()} function applies an orthographic rotation defined by
azimuth and elevation angles, returning a 2D projection matrix for custom
plotting:

\begin{verbatim}
edges <- edges.torus(8, 12)
coords <- grip.layout(edges, n = max(edges), dim = 3,
                       preset = "torus", seed = 3)
grip.plot(coords, edges, projection = "ortho", main = "Torus (8x12)")
\end{verbatim}

\begin{figure}

{\centering \includegraphics[width=1\linewidth]{grip-paper_20260326_092635_files/figure-latex/torus-plot-1} 

}

\caption{An 8x12 torus laid out in 3D with the torus preset, shown as a static orthographic projection.}(\#fig:torus-plot)
\end{figure}

\subsection{Graph generators}\label{graph-generators}

The package includes convenience functions for generating common graph
families as two-column edge matrices: \texttt{edges.mesh()} (rectangular grid),
\texttt{edges.cycle()}, \texttt{edges.path()}, \texttt{edges.kary.tree()} (k-ary tree),
\texttt{edges.torus()}, \texttt{edges.cylinder()}, \texttt{edges.sierpinski.triangle()},
\texttt{edges.sierpinski.tetrahedron()}, and \texttt{edges.sierpinski.carpet()}. These are
useful for benchmarking, testing, and examples.

\section{Layout quality metrics}\label{sec:metrics}

For real-world graphs---social networks, biological networks, knowledge
graphs---there is typically no ground-truth embedding. Layout quality must be
assessed through proxy metrics that capture different aspects of a good
drawing. The \textbf{grip} package provides six such metrics, accessible through
\texttt{grip.score.layout()} for individual layouts and aggregated by
\texttt{grip.compare.layouts()} for systematic comparisons.

\subsection{Sampled stress}\label{sampled-stress}

Stress measures how faithfully the Euclidean distances in the layout
reproduce the graph-theoretic distances. For scalability, \textbf{grip} computes a
sampled approximation: a random subset of vertex pairs is drawn, shortest-path
distances are computed via BFS (unweighted) or Dijkstra's algorithm
(weighted), and the normalized root-mean-square discrepancy between the
(optimally scaled) Euclidean distances and graph distances is reported. Lower
values indicate better fidelity. The sample size is controlled by
\texttt{sample.size.stress} (default 2000 pairs).

\subsection{Edge-length coefficient of variation}\label{edge-length-coefficient-of-variation}

This metric measures the uniformity of edge lengths in the layout. The
coefficient of variation (standard deviation divided by the mean) of all edge
lengths is computed. Lower values indicate that neighboring vertices are
drawn with more consistent spacing, which generally improves readability.

\subsection{Sampled non-edge separation ratio}\label{sampled-non-edge-separation-ratio}

For a layout to be readable, non-neighboring vertices should not be drawn
too close together. This metric samples random non-edge pairs, computes
their Euclidean distances, and reports the ratio of the minimum sampled
non-edge distance to the median edge length. Higher values indicate more
breathing room between non-neighbors.

\subsection{Edge crossings}\label{edge-crossings}

In 2D layouts, edge crossings reduce readability. \textbf{grip} counts exact edge
crossings using a pairwise segment intersection test. Because this is
\(O(m^2)\) in the number of edges, it is computed automatically only for
small graphs (controlled by \texttt{edge.crossings.max.edges}; default 1000 edges)
and can be disabled entirely with \texttt{edge.crossings\ =\ "never"}.

\subsection{Cluster separation}\label{cluster-separation}

When external labels are available (e.g., community memberships, biological
classifications), cluster separation measures how well the layout separates
known groups. It is defined as the ratio of the mean pairwise distance
between cluster centroids to the mean within-cluster spread. Higher values
indicate clearer visual separation of groups. This metric is reported only
when \texttt{clusters} are supplied.

\subsection{Procrustes stability}\label{procrustes-stability}

A good layout should be reproducible: running the algorithm with different
random seeds should yield geometrically similar results. Procrustes stability
is assessed by computing layouts across multiple seeds, optimally aligning
each pair using Procrustes analysis (rotation, reflection, and uniform
scaling), and reporting the mean RMSE of the aligned coordinates. Lower
values indicate greater stability.

\subsection{Composite score}\label{composite-score}

The \texttt{grip.compare.layouts()} function computes a rank-based composite score
across all available metrics, using configurable weights. The default weights
emphasize stress (0.35), edge-length CV (0.20), non-edge separation (0.20),
and Procrustes stability (0.15), with smaller weights for edge crossings
(0.05) and cluster separation (0.05). The composite score is useful as a
screening tool for shortlisting candidates but should not be treated as the
sole basis for layout selection.

\subsection{Scoring API}\label{scoring-api}

Individual layouts are scored with \texttt{grip.score.layout()}:

\begin{verbatim}
edges <- edges.mesh(5, 5)
coords <- grip.layout(edges, n = 25, dim = 2, preset = "mesh", seed = 1)
score <- grip.score.layout(coords, edges = edges, n = 25)
paper_kable(
  score,
  caption = paper_table_caption(
    "Quality metrics for a single 5x5 mesh layout.",
    "nV = vertices; nE = edges; Dim = layout dimension; Stress = sampled stress; EdgeCV = edge-length coefficient of variation; MedEL = median edge length; NESR = sampled non-edge separation ratio."
  ),
  digits = 3
)
\end{verbatim}

\begin{table}

\caption{(\#tab:unnamed-chunk-3)Quality metrics for a single 5x5 mesh layout. Abbrev.: nV = vertices; nE = edges; Dim = layout dimension; Stress = sampled stress; EdgeCV = edge-length coefficient of variation; MedEL = median edge length; NESR = sampled non-edge separation ratio.}
\centering
\begin{tabular}[t]{l|r|r|r|r|r|r|r|r}
\hline
nV & nE & Dim & Stress & EdgeCV & MedEL & NESR & Cross & ClSep\\
\hline
25 & 40 & 2 & 9.192 & 0.055 & 26.52 & 1.286 & 0 & NA\\
\hline
\end{tabular}
\end{table}

Multiple candidates are compared with \texttt{grip.compare.layouts()}, which runs
each candidate across multiple seeds, computes per-run scores, and aggregates
them into a summary table with stability metrics:

\begin{verbatim}
edges <- edges.mesh(6, 6)
cmp <- grip.compare.layouts(
  edges = edges, n = 36, dim = 2,
  candidates = c("default", "mesh"),
  seeds = 1:3
)
paper_kable(
  cmp$summary[, c("candidate", "score.composite",
                  "sampled.stress.mean", "edge.length.cv.mean",
                  "sampled.nonedge.sep.ratio.mean",
                  "stability.procrustes.mean")],
  caption = paper_table_caption(
    "Comparison of default and mesh candidates across seeds on a 6x6 mesh.",
    "Cand = candidate; Comp = composite score; Stress = sampled stress; EdgeCV = edge-length coefficient of variation; NESR = sampled non-edge separation ratio; PStab = Procrustes stability."
  ),
  digits = 3
)
\end{verbatim}

\begin{table}

\caption{(\#tab:unnamed-chunk-4)Comparison of default and mesh candidates across seeds on a 6x6 mesh. Abbrev.: Cand = candidate; Comp = composite score; Stress = sampled stress; EdgeCV = edge-length coefficient of variation; NESR = sampled non-edge separation ratio; PStab = Procrustes stability.}
\centering
\begin{tabular}[t]{l|r|r|r|r|r}
\hline
Cand & Comp & Stress & EdgeCV & NESR & PStab\\
\hline
mesh & 0 & 10.47 & 0.055 & 1.237 & 0.002\\
\hline
default & 1 & 17.69 & 0.181 & 0.168 & 0.552\\
\hline
\end{tabular}
\end{table}

\section{Worked examples}\label{worked-examples}

\subsection{Synthetic example: default versus preset}\label{synthetic-example-default-versus-preset}

To illustrate the value of presets, we compare the default GRIP parameters
against the carpet preset on a level-3 Sierpinski carpet (512 vertices, 704
edges). The carpet's recursive structure with large interior holes is
challenging for generic parameters.

\begin{verbatim}
edges <- edges.sierpinski.carpet(3)
n <- max(edges)

carpet.default <- grip.layout(edges, n = n, dim = 2, seed = 12)
carpet.preset  <- grip.layout(edges, n = n, dim = 2, preset = "carpet", seed = 12)

op <- par(mfrow = c(1, 2), mar = c(1, 1, 3, 1))
grip.plot(carpet.default, edges,
          main = "Carpet default", pch = 16, cex = 0.3,
          axes = FALSE, xlab = "", ylab = "", frame.plot = FALSE)
grip.plot(carpet.preset, edges,
          main = "Carpet preset", pch = 16, cex = 0.3,
          axes = FALSE, xlab = "", ylab = "", frame.plot = FALSE)
\end{verbatim}

\begin{figure}

{\centering \includegraphics[width=1\linewidth]{grip-paper_20260326_092635_files/figure-latex/carpet-compare-1} 

}

\caption{Level-3 Sierpinski carpet with default parameters (left) and the carpet preset (right). The preset recovers the fractal hole structure more clearly.}(\#fig:carpet-compare)
\end{figure}

\begin{verbatim}
par(op)
\end{verbatim}

The quality metrics confirm the visual impression:

\begin{verbatim}
carpet.cmp <- grip.compare.layouts(
  edges = edges, n = n, dim = 2,
  candidates = c("default", "carpet"),
  seeds = 1:3
)
paper_kable(
  carpet.cmp$summary[, c("candidate", "score.composite",
                          "sampled.stress.mean", "edge.length.cv.mean",
                          "sampled.nonedge.sep.ratio.mean",
                          "stability.procrustes.mean")],
  caption = paper_table_caption(
    "Comparison of default and carpet presets on the level-3 Sierpinski carpet.",
    "Cand = candidate; Comp = composite score; Stress = sampled stress; EdgeCV = edge-length coefficient of variation; NESR = sampled non-edge separation ratio; PStab = Procrustes stability."
  ),
  digits = 3
)
\end{verbatim}

\begin{table}

\caption{(\#tab:unnamed-chunk-5)Comparison of default and carpet presets on the level-3 Sierpinski carpet. Abbrev.: Cand = candidate; Comp = composite score; Stress = sampled stress; EdgeCV = edge-length coefficient of variation; NESR = sampled non-edge separation ratio; PStab = Procrustes stability.}
\centering
\begin{tabular}[t]{l|r|r|r|r|r}
\hline
Cand & Comp & Stress & EdgeCV & NESR & PStab\\
\hline
carpet & 0 & 54.677 & 0.119 & 1.092 & 0.022\\
\hline
default & 1 & 121.332 & 0.191 & 0.108 & 0.537\\
\hline
\end{tabular}
\end{table}

\subsection{Real-data example: Zachary karate club}\label{sec:karate}

The Zachary karate club network \citep{zachary1977} is a 34-vertex social network
with a known split into two factions. This example demonstrates the full
comparison pipeline on a graph where the layout can be validated against
meaningful external metadata.

\begin{verbatim}
find.extdata.file <- function(...) {
  rel.path <- file.path(...)
  candidates <- c(
    system.file("extdata", rel.path, package = "grip"),
    file.path("..", "inst", "extdata", rel.path),
    file.path("inst", "extdata", rel.path)
  )
  path <- candidates[file.exists(candidates)][1L]
  if (!length(path) || is.na(path) || !nzchar(path))
    stop("could not locate ", rel.path)
  path
}

read.extdata.csv <- function(file.name) {
  utils::read.csv(find.extdata.file(file.name), stringsAsFactors = FALSE)
}

read.extdata.rds <- function(...) {
  readRDS(find.extdata.file(...))
}

karate.edges <- as.matrix(read.extdata.csv("karate-club-edges.csv"))
karate.n <- max(karate.edges)
labels.df <- read.extdata.csv("karate-club-membership.csv")
karate.club <- labels.df$club[order(labels.df$vertex)]
\end{verbatim}

\subsubsection{Step 1: shortlist presets}\label{step-1-shortlist-presets}

We compare the default layout against the tree and mesh presets in 3D:

\begin{verbatim}
karate.presets <- grip.compare.layouts(
  edges = karate.edges, n = karate.n, dim = 3,
  candidates = c("default", "tree", "mesh"),
  seeds = 1:3,
  sample.size.stress = 1000L,
  sample.size.nonedge = 2000L,
  edge.crossings = "never",
  return.layouts = TRUE
)
paper_kable(
  karate.presets$summary[, c("candidate", "score.composite",
    "sampled.stress.mean", "edge.length.cv.mean",
    "sampled.nonedge.sep.ratio.mean",
    "stability.procrustes.mean")],
  caption = paper_table_caption(
    "Preset shortlisting on the Zachary karate club graph in 3D.",
    "Cand = candidate; Comp = composite score; Stress = sampled stress; EdgeCV = edge-length coefficient of variation; NESR = sampled non-edge separation ratio; PStab = Procrustes stability."
  ),
  digits = 3
)
\end{verbatim}

\begin{table}

\caption{(\#tab:unnamed-chunk-7)Preset shortlisting on the Zachary karate club graph in 3D. Abbrev.: Cand = candidate; Comp = composite score; Stress = sampled stress; EdgeCV = edge-length coefficient of variation; NESR = sampled non-edge separation ratio; PStab = Procrustes stability.}
\centering
\begin{tabular}[t]{l|r|r|r|r|r}
\hline
Cand & Comp & Stress & EdgeCV & NESR & PStab\\
\hline
tree & 0.333 & 0.759 & 0.492 & 0.004 & 0.073\\
\hline
default & 0.556 & 11.411 & 0.207 & 0.206 & 0.427\\
\hline
mesh & 0.611 & 10.664 & 0.227 & 0.000 & 0.216\\
\hline
\end{tabular}
\end{table}

\subsubsection{Step 2: local parameter search}\label{step-2-local-parameter-search}

Starting from the top-ranked preset, we search a small parameter neighborhood:

\begin{verbatim}
karate.base <- karate.presets$summary[1, , drop = FALSE]
karate.search <- list(
  candidate.prefix = "karate.local",
  placement = karate.base$placement[[1L]],
  rounds = karate.base$rounds[[1L]],
  final_rounds = sort(unique(as.integer(c(
    max(32L, karate.base$final.rounds[[1L]] - 32L),
    karate.base$final.rounds[[1L]],
    karate.base$final.rounds[[1L]] + 32L
  )))),
  num_init = karate.base$num.init[[1L]],
  num_nbrs = sort(unique(as.integer(c(
    max(4L, karate.base$num.nbrs[[1L]] - 4L),
    karate.base$num.nbrs[[1L]],
    karate.base$num.nbrs[[1L]] + 4L
  )))),
  r = karate.base$r[[1L]],
  s = karate.base$s[[1L]],
  repulsion_factor = sort(unique(c(
    max(0, karate.base$repulsion.factor[[1L]] - 0.5),
    karate.base$repulsion.factor[[1L]],
    karate.base$repulsion.factor[[1L]] + 0.5
  )))
)
if (!is.na(karate.base$preset[[1L]]) && nzchar(karate.base$preset[[1L]]))
  karate.search$preset <- karate.base$preset[[1L]]

karate.local <- grip.compare.layouts(
  edges = karate.edges, n = karate.n, dim = 3,
  search = karate.search,
  seeds = 1:2,
  sample.size.stress = 1000L,
  sample.size.nonedge = 2000L,
  edge.crossings = "never",
  return.layouts = TRUE
)
karate.local.display <- head(karate.local$summary[, c(
  "candidate", "score.composite",
  "sampled.stress.mean", "sampled.nonedge.sep.ratio.mean",
  "stability.procrustes.mean"
)], 6)
karate.local.display$candidate <- compact_search_candidate(karate.local.display$candidate)

paper_kable(
  karate.local.display,
  caption = paper_table_caption(
    "Top local-search candidates around the leading karate preset.",
    "Cand = candidate code; fr = final rounds; nn = neighborhood count; rf = repulsion factor; Comp = composite score; Stress = sampled stress; NESR = sampled non-edge separation ratio; PStab = Procrustes stability."
  ),
  digits = 3
)
\end{verbatim}

\begin{table}

\caption{(\#tab:unnamed-chunk-8)Top local-search candidates around the leading karate preset. Abbrev.: Cand = candidate code; fr = final rounds; nn = neighborhood count; rf = repulsion factor; Comp = composite score; Stress = sampled stress; NESR = sampled non-edge separation ratio; PStab = Procrustes stability.}
\centering
\begin{tabular}[t]{l|r|r|r|r}
\hline
Cand & Comp & Stress & NESR & PStab\\
\hline
fr192-nn4-rf0 & 0.281 & 0.685 & 0.007 & 0.046\\
\hline
fr160-nn4-rf0 & 0.330 & 0.752 & 0.008 & 0.046\\
\hline
fr192-nn12-rf0 & 0.366 & 0.698 & 0.003 & 0.051\\
\hline
fr192-nn8-rf0 & 0.366 & 0.719 & 0.004 & 0.078\\
\hline
fr160-nn8-rf05 & 0.382 & 7.683 & 0.037 & 0.262\\
\hline
fr128-nn4-rf0 & 0.408 & 0.841 & 0.008 & 0.046\\
\hline
\end{tabular}
\end{table}

\subsubsection{Step 3: incorporate external metadata}\label{step-3-incorporate-external-metadata}

The karate data include faction labels. We re-compare the best local-search
candidate against the default, now including \texttt{cluster.separation}:

\begin{verbatim}
karate.local.best <- karate.local$summary[1, , drop = FALSE]
karate.metadata <- grip.compare.layouts(
  edges = karate.edges, n = karate.n, dim = 3,
  candidates = list(
    default = NULL,
    tuned = grip.params.from.summary(karate.local.best)
  ),
  clusters = karate.club,
  seeds = 1:3,
  sample.size.stress = 1000L,
  sample.size.nonedge = 2000L,
  edge.crossings = "never",
  return.layouts = TRUE
)
paper_kable(
  karate.metadata$summary[, c("candidate", "score.composite",
    "sampled.stress.mean", "cluster.separation.mean")],
  caption = paper_table_caption(
    "Metadata-aware comparison of the default and tuned karate layouts.",
    "Cand = candidate; Comp = composite score; Stress = sampled stress; ClSep = cluster separation."
  ),
  digits = 3
)
\end{verbatim}

\begin{table}

\caption{(\#tab:unnamed-chunk-9)Metadata-aware comparison of the default and tuned karate layouts. Abbrev.: Cand = candidate; Comp = composite score; Stress = sampled stress; ClSep = cluster separation.}
\centering
\begin{tabular}[t]{l|r|r|r}
\hline
Cand & Comp & Stress & ClSep\\
\hline
tuned & 0.421 & 0.603 & 3.652\\
\hline
default & 0.579 & 11.411 & 1.134\\
\hline
\end{tabular}
\end{table}

The metadata comparison above keeps compact tree-derived layouts on the
shortlist, but for a static orthographic figure a slightly more repulsive
tree-based variant opens the dense core and preserves the faction split more
clearly. We therefore render the tuned candidate below and choose the seed
with the strongest cluster separation among three runs.

\begin{verbatim}
karate.figure.params <- list(
  preset = "tree",
  repulsion_factor = 4,
  rounds = 160L,
  final_rounds = 224L,
  num_nbrs = 16L,
  r = 0.08,
  s = 6
)

karate.figure <- grip.compare.layouts(
  edges = karate.edges, n = karate.n, dim = 3,
  candidates = list(tuned = karate.figure.params),
  clusters = karate.club,
  seeds = 1:3,
  sample.size.stress = 1000L,
  sample.size.nonedge = 2000L,
  edge.crossings = "never",
  return.layouts = TRUE
)

ok.runs <- karate.figure$runs[karate.figure$runs$status == "ok", , drop = FALSE]
best.run <- ok.runs[order(
  -ok.runs$cluster.separation,
  -ok.runs$sampled.nonedge.sep.ratio,
  ok.runs$sampled.stress,
  ok.runs$seed
), , drop = FALSE][1, , drop = FALSE]

best.coords <- karate.figure$layouts$tuned[[as.character(best.run$seed[[1L]])]]
best.coords <- prcomp(best.coords, center = TRUE, scale. = FALSE)$x
club.cols <- ifelse(karate.club == "Mr. Hi", "#1b9e77", "#d95f02")

op <- par(mfrow = c(1, 2), mar = c(1, 1, 3, 1))
grip.plot(best.coords, karate.edges, projection = "ortho",
          azimuth = 25, elevation = 20,
          vertex.col = club.cols, main = "View 1")
grip.plot(best.coords, karate.edges, projection = "ortho",
          azimuth = -55, elevation = 30,
          vertex.col = club.cols, main = "View 2")
\end{verbatim}

\begin{figure}

{\centering \includegraphics[width=1\linewidth]{grip-paper_20260326_092635_files/figure-latex/karate-final-1} 

}

\caption{Zachary karate club with a tuned tree-derived 3D layout, colored by faction membership (green = Mr.\ Hi, orange = Officer). Two orthographic views of the static layout.}(\#fig:karate-final)
\end{figure}

\begin{verbatim}
par(op)
\end{verbatim}

\subsection{Scaled example: HMP microbial network}\label{sec:hmp}

The package bundles a coarsened real-world graph derived from the Human
Microbiome Project (HMP) 16S amplicon analysis. The original giant connected
component contained 6,474 vertices from an iterative k-nearest-neighbor (iKNN)
graph with \(k=3\) built on a \(\geq 1\%\) relative-abundance OTU table. Greedy
graph coarsening in two rounds reduced this to 1,828 supernodes, preserving
the large-scale topology including arm-like extensions corresponding to
distinct community state types (CSTs).

This dataset is large enough that interactive manual layout selection is
impractical, making \textbf{grip}'s systematic comparison workflow especially
valuable.

\begin{verbatim}
data(hmp.u01.gc.coarse)
cat("Vertices:", length(hmp.u01.gc.coarse$adj_list), "\n")
#> Vertices: 1828
cat("CST levels:", paste(sort(unique(hmp.u01.gc.coarse$vertex_data$cst)),
                          collapse = ", "), "\n")
#> CST levels: I, II, III, IV-A, IV-B, IV-C, V
\end{verbatim}

The full comparison pipeline (preset shortlisting followed by local search)
takes several minutes on this graph, so the package bundles precomputed
results for the vignette. The code to reproduce the comparison is shown with
\texttt{eval=FALSE} and is fully documented in the package's HMP vignette.

\begin{verbatim}
# Preset shortlisting (not run in paper build)
cmp <- grip.compare.layouts(
  adj_list = hmp.u01.gc.coarse$adj_list,
  weight_list = hmp.u01.gc.coarse$weight_list,
  n = length(hmp.u01.gc.coarse$adj_list),
  dim = 3,
  candidates = c("default", "tree", "torus"),
  clusters = hmp.u01.gc.coarse$vertex_data$cst,
  seeds = 1,
  sample.size.stress = 1200,
  sample.size.nonedge = 3000
)
\end{verbatim}

\begin{verbatim}
# Local search around strongest region (not run in paper build)
search.cmp <- grip.compare.layouts(
  adj_list = hmp.u01.gc.coarse$adj_list,
  weight_list = hmp.u01.gc.coarse$weight_list,
  n = length(hmp.u01.gc.coarse$adj_list),
  dim = 3,
  search = list(
    candidate.prefix = "hmp.local",
    placement = "barycenter",
    rounds = 192L,
    final_rounds = c(224L, 288L),
    num_init = 30L,
    num_nbrs = 10L,
    r = 0.10, s = 2.0,
    repulsion_factor = c(0.75, 1.25)
  ),
  clusters = hmp.u01.gc.coarse$vertex_data$cst,
  seeds = 1,
  sample.size.stress = 1200,
  sample.size.nonedge = 3000
)
\end{verbatim}

The precomputed results show that the tree and torus presets both improve on
the default for this graph, and that moderate repulsion (rather than maximum
repulsion) produces the best balance between stress and CST cluster
separation. The final layout colored by CST labels shows clear arm-like
structures corresponding to distinct microbial community types.

\begin{verbatim}
hmp.results <- tryCatch({
  path <- system.file("extdata", "hmp_u01_gc_coarse",
                       "vignette_results.rds", package = "grip")
  if (nzchar(path)) readRDS(path) else NULL
}, error = function(e) NULL)

if (!is.null(hmp.results)) {
  edge.matrix.from.adj <- function(adj.list) {
    edges <- list(); idx <- 0L
    for (u in seq_along(adj.list)) {
      nbrs <- adj.list[[u]][adj.list[[u]] > u]
      for (v in nbrs) { idx <- idx + 1L; edges[[idx]] <- c(u, v) }
    }
    do.call(rbind, edges)
  }
  hmp.edges <- edge.matrix.from.adj(hmp.u01.gc.coarse$adj_list)
  hmp.cst <- hmp.u01.gc.coarse$vertex_data$cst
  cst.levels <- sort(unique(hmp.cst))
  cst.cols <- setNames(
    grDevices::hcl.colors(length(cst.levels), "Dark 3"),
    cst.levels
  )
  hmp.cst.col <- cst.cols[hmp.cst]

  op <- par(mfrow = c(1, 3), mar = c(1, 1, 3, 1))
  grip.plot(hmp.results$layouts$preset$default, hmp.edges,
            projection = "ortho", vertex.col = hmp.cst.col, main = "default")
  grip.plot(hmp.results$layouts$preset$tree, hmp.edges,
            projection = "ortho", vertex.col = hmp.cst.col, main = "tree")
  grip.plot(hmp.results$layouts$preset$torus, hmp.edges,
            projection = "ortho", vertex.col = hmp.cst.col, main = "torus")
  par(op)
} else {
  message("HMP precomputed results not available; skipping figure.")
}
\end{verbatim}

\begin{figure}

{\centering \includegraphics[width=1\linewidth]{grip-paper_20260326_092635_files/figure-latex/hmp-layouts-1} 

}

\caption{Precomputed 3D layouts of the 1,828-vertex HMP microbial network colored by community state type (CST). Left: default parameters. Center: tree preset. Right: torus preset. The tree and torus presets both produce cleaner separation of the CST-associated arms.}(\#fig:hmp-layouts)
\end{figure}

\section{Comparison with existing tools}\label{sec:comparison}

To provide context for \textbf{grip}'s performance, we compare it against three
layout algorithms available in \textbf{igraph} \citep{csardi2006}: Fruchterman--Reingold
(\texttt{layout\_with\_fr}), Kamada--Kawai (\texttt{layout\_with\_kk}), and DrL
(\texttt{layout\_with\_drl}). We use \texttt{grip.score.layout()} to evaluate all layouts
on the same quality metrics, ensuring a fair comparison.

The comparison uses the Zachary karate club graph (34 vertices, 78 edges)
in 2D. Note that \textbf{grip}'s quality metrics can evaluate any layout, not
just those produced by \textbf{grip}.

\begin{verbatim}
benchmark.results <- read.extdata.rds("vs_alternatives", "benchmark_results.rds")

benchmark.methods <- c(
  "FR (igraph)",
  "KK (igraph)",
  "DrL (igraph)",
  "grip default",
  "grip tuned"
)
score.all <- benchmark.results$karate$scores[
  match(benchmark.methods, benchmark.results$karate$scores$method),
  c("method", "sampled.stress", "edge.length.cv",
    "sampled.nonedge.sep.ratio", "cluster.separation",
    "edge.crossings")
]
score.all$method <- c(
  "FR (igraph)",
  "KK (igraph)",
  "DrL (igraph)",
  "GRIP default",
  "GRIP tuned"
)

paper_kable(
  score.all,
  caption = paper_table_caption(
    "Illustrative 2D comparison on the Zachary karate club graph.",
    "Meth = method; Stress = sampled stress; EdgeCV = edge-length coefficient of variation; NESR = sampled non-edge separation ratio; ClSep = cluster separation; Cross = edge crossings."
  ),
  digits = 3
)
\end{verbatim}

\begin{table}

\caption{(\#tab:unnamed-chunk-13)Illustrative 2D comparison on the Zachary karate club graph. Abbrev.: Meth = method; Stress = sampled stress; EdgeCV = edge-length coefficient of variation; NESR = sampled non-edge separation ratio; ClSep = cluster separation; Cross = edge crossings.}
\centering
\begin{tabular}[t]{l|r|r|r|r|r}
\hline
Meth & Stress & EdgeCV & NESR & ClSep & Cross\\
\hline
FR (igraph) & 1.220 & 0.358 & 0.367 & 2.757 & 68\\
\hline
KK (igraph) & 0.614 & 0.281 & 0.302 & 2.253 & 73\\
\hline
DrL (igraph) & 0.605 & 0.468 & 0.237 & 1.658 & 120\\
\hline
GRIP default & 13.825 & 0.279 & 0.083 & 0.266 & 301\\
\hline
GRIP tuned & 18.123 & 0.285 & 0.118 & 0.561 & 246\\
\hline
\end{tabular}
\end{table}

On this very small 2D social network, the classical igraph layouts are
stronger on stress and edge crossings, which is consistent with their design
goals. The tuned \textbf{grip} configuration improves on the package defaults for
cluster separation and crossings, but it does not outperform FR, KK, or DrL
on this benchmark. This is an informative result rather than a negative one:
\textbf{grip} is aimed at larger graphs, structured graph families, and 3D layout
selection workflows. Its key advantage here is that \texttt{grip.score.layout()}
provides a common scoring framework that can assess layouts from any source,
making \textbf{grip} complementary to existing tools rather than a replacement for
them.

\section{Discussion and future work}\label{discussion-and-future-work}

The \textbf{grip} package provides an accessible R implementation of the GRIP
multiscale force-directed layout algorithm, complemented by validated
parameter presets and a principled quality assessment framework. The
combination of layout computation and quality scoring in a single package
addresses a practical gap in the R ecosystem: for real-world graphs where no
canonical embedding exists, choosing among layout options has historically
been a visual and subjective process.

Several limitations should be noted. The current set of four presets covers
a finite range of graph topologies; real-world graphs may not resemble any of
these families closely. The quality metrics are heuristic proxies for layout
quality, not formal optimality guarantees. Very large graphs (100,000+
vertices) may still require substantial computation time, as the current
implementation does not exploit GPU parallelism. The edge crossing count is
exact but scales quadratically in the number of edges, limiting its practical
use to moderately sized graphs.

Future development directions include additional presets (e.g., for
scale-free networks and bipartite graphs), integration with Bioconductor
graph classes for seamless use in genomics workflows, GPU-accelerated force
computation for very large graphs, additional quality metrics such as angular
resolution and neighborhood preservation, and potential automatic preset
selection based on graph-structural features. A Shiny-based interactive
explorer (\textbf{gripui}) is available as a companion package but is not the
focus of this article.

The package source code is available at
\url{https://github.com/pgajer/grip} and is being prepared for CRAN
submission.

\section{Acknowledgments}\label{acknowledgments}

The GRIP algorithm was originally developed in collaboration with Stephen G.
Kobourov and Michael T. Goodrich. The HMP microbial network data were derived
from publicly available Human Microbiome Project resources. Research reported
in this publication was supported by grant number INV-048956 from the Gates
Foundation.

\bibliography{grip-paper.bib}

\address{%
Pawel Gajer\\
Center for Advanced Microbiome Research and Innovation (CAMRI), Institute for Genome Sciences, Department of Microbiology and Immunology, University of Maryland School of Medicine\\%
Baltimore, Maryland, USA\\
%
%
%
\href{mailto:pgajer@gmail.com}{\nolinkurl{pgajer@gmail.com}}%
}

\address{%
Jacques Ravel\\
Center for Advanced Microbiome Research and Innovation (CAMRI), Institute for Genome Sciences, Department of Microbiology and Immunology, University of Maryland School of Medicine\\%
Baltimore, Maryland, USA\\
%
%
%
%
}
